% =====================================================================
%  CHAPTER 9: 食品添加剂及其管理 (Food Additives and Management)
% =====================================================================
\chapter{食品添加剂及其管理}
\setcounter{chapter}{9}
\chapterintro{
\keyword{掌握}食品添加剂的定义、使用要求；掌握JECFA安全性评价分类；掌握各类食品添加剂（抗氧化剂、防腐剂、护色剂——亚硝酸盐*等）的作用。\keyword{熟悉}食品添加剂的管理制度。\keyword{了解}各类食品添加剂的常用品种和使用范围。
}

% =====================================================================
%  第一节 食品添加剂的定义与分类
% =====================================================================
\section{食品添加剂的定义与分类}

\subsection{定义}

\begin{keybox}
\textbf{食品添加剂的定义（掌握）——中国《食品安全法》：}
食品添加剂是指为改善食品品质和色、香、味，以及为防腐、保鲜和加工工艺的需要而加入食品中的\textbf{人工合成或者天然物质}。

\textbf{定义涵盖范围包括：}
\begin{itemize}[leftmargin=*]
    \item 食品用香料
    \item 胶基糖果中基础剂物质
    \item 食品工业用加工助剂
    \item 营养强化剂
\end{itemize}
\end{keybox}

\begin{tipbox}
\textbf{中国与CAC定义的差异：}
\begin{itemize}[leftmargin=*]
    \item \textbf{中国定义：}包含营养强化剂、食品用香料、加工助剂、胶基物质。
    \item \textbf{CAC（国际食品法典委员会）定义：}食品添加剂通常不包括污染物，也不包括为保持或提高营养价值而添加的物质（即营养强化剂不属于添加剂范畴）。
    \item 中国的定义范围更广，将营养强化剂也包括在食品添加剂中。
\end{itemize}
\end{tipbox}

\subsection{分类}

\textbf{中国GB 2760-2014《食品安全国家标准 食品添加剂使用标准》将食品添加剂按功能分为23类，主要包括：}

\begin{enumerate}[leftmargin=*]
    \item 酸度调节剂（Acidity Regulator）
    \item 抗结剂（Anticaking Agent）
    \item 消泡剂（Antifoaming Agent）
    \item 抗氧化剂（Antioxidant）
    \item 漂白剂（Bleaching Agent）
    \item 膨松剂（Bulking Agent）
    \item 着色剂（Color）
    \item 护色剂（Color Fixative）
    \item 乳化剂（Emulsifier）
    \item 酶制剂（Enzyme Preparation）
    \item 增味剂（Flavor Enhancer）
    \item 面粉处理剂（Flour Treatment Agent）
    \item 被膜剂（Coating Agent）
    \item 水分保持剂（Humectant）
    \item 营养强化剂（Nutrient Fortifier）
    \item 防腐剂（Preservative）
    \item 稳定剂和凝固剂（Stabilizer and Coagulator）
    \item 甜味剂（Sweetener）
    \item 增稠剂（Thickener）
    \item 食品用香料（Food Flavorings）
    \item 食品工业用加工助剂（Food Processing Aids）
    \item 胶基糖果中基础剂物质（Gum Base Materials）
    \item 其他（Others）
\end{enumerate}

% =====================================================================
%  第二节 食品添加剂的使用要求
% =====================================================================
\section{食品添加剂的使用要求}

\subsection{基本要求}

\begin{keybox}
\textbf{食品添加剂的使用原则（掌握）：}
\begin{enumerate}[leftmargin=*]
    \item \textbf{不应对人体产生任何健康危害。}
    \item \textbf{不应掩盖食品腐败变质。}
    \item \textbf{不应掩盖食品本身或加工过程中的质量缺陷（即不得以掺杂、掺假、伪造为目的使用食品添加剂！）。}
    \item \textbf{不应降低食品本身的营养价值。}
    \item \textbf{在达到预期效果的前提下，尽可能降低在食品中的使用量。}
\end{enumerate}
\end{keybox}

\begin{keybox}
\textbf{食品添加剂允许使用的情形（掌握）：}
\begin{enumerate}[leftmargin=*]
    \item 保持或提高食品本身的营养价值
    \item 作为某些特殊膳食用食品的必需配料或成分
    \item 提高食品的质量和稳定性，改进其感官特性
    \item 便于食品的生产、加工、包装、运输或者贮藏
\end{enumerate}
\end{keybox}

\subsection{带入原则（Carry-over Principle）}

\begin{defbox}
\textbf{带入原则（Carry-over Principle）：}某种食品添加剂不是直接加入到某食品中，而是通过含该添加剂的食品配料（或原料）带入到该食品中的情况。

\textbf{带入原则的适用条件：}
\begin{enumerate}[leftmargin=*]
    \item 根据GB 2760，食品配料中允许使用该添加剂
    \item 食品配料中该添加剂的用量不应超过允许的最大使用量
    \item 应在正常生产工艺条件下使用这些配料
    \item 由配料带入食品中的该添加剂的含量应明显低于直接将其添加到该食品中通常所需要的水平
\end{enumerate}
\end{defbox}

% =====================================================================
%  第三节 JECFA安全性评价分类
% =====================================================================
\section{JECFA对食品添加剂的安全性评价分类}

\begin{keybox}
\textbf{JECFA（FAO/WHO联合食品添加剂专家委员会）安全性评价分类（掌握）*：}
\end{keybox}

\begin{enumerate}[leftmargin=*]
    \item \textbf{第一类 GRAS（Generally Recognized As Safe，一般公认安全）：}
    \begin{itemize}[leftmargin=*]
        \item 以现有毒理学资料和化学、生物化学知识来看，按正常需要使用，对人体是安全的。
        \item 此类物质毒性极低、不规定ADI值。
        \item 可按照正常生产需要使用，不需限定最大用量。
        \item 示例：食盐、糖、柠檬酸、醋酸、碳酸氢钠等。
    \end{itemize}

    \item \textbf{第二类 A类：}
    \begin{itemize}[leftmargin=*]
        \item \textbf{A1类：}JECFA已制定正式ADI值（有完整的毒理学资料，安全性评价明确）。
        \item \textbf{A2类：}JECFA尚未制定正式ADI值（毒理学资料不完全），仅制定临时ADI值，在规定的期限内暂时允许使用。
    \end{itemize}

    \item \textbf{第三类 B类：}
    \begin{itemize}[leftmargin=*]
        \item JECFA曾进行过安全性评价，但因毒理学资料不足，\textbf{未能制定ADI值}。
        \item 需要进一步的毒理学数据补充后才能重新评价。
    \end{itemize}

    \item \textbf{第四类 C类：}
    \begin{itemize}[leftmargin=*]
        \item \textbf{C1类：}JECFA根据毒理学资料，认为在食品中使用是\textbf{不安全的}。
        \item \textbf{C2类：}JECFA根据毒理学资料，认为应\textbf{仅限于特殊用途食品或在特殊条件下}严格限制使用。
    \end{itemize}
\end{enumerate}

\begin{exambox}
\textbf{常见考试混淆点：}
\begin{itemize}[leftmargin=*]
    \item GRAS（第一类）= 无ADI（不需要限制），正常使用即可
    \item A1 = 有正式ADI；A2 = 仅临时ADI
    \item B类 = 毒理学资料不足，无ADI
    \item C1 = 不安全；C2 = 仅特殊情况限制使用
\end{itemize}
\end{exambox}

% =====================================================================
%  第四节 各类食品添加剂（常用品种）
% =====================================================================
\section{常用食品添加剂各论}

\subsection{酸度调节剂（Acidity Regulator）}

\begin{itemize}[leftmargin=*]
    \item \textbf{常用品种：}柠檬酸、乳酸、苹果酸、酒石酸、磷酸（酸化剂）；碳酸氢钠、碳酸钠、氢氧化钠（碱化剂）。
    \item \textbf{特点：}大多数酸度调节剂为机体正常代谢产物或参与正常代谢的物质，\textbf{毒性极低}，使用安全性高。
    \item \textbf{功能：}调节/维持食品酸度，改善风味，凝胶作用，防腐辅助（降低pH可抑制某些微生物）。
\end{itemize}

\subsection{抗氧化剂（Antioxidant）}

\begin{keybox}
\textbf{抗氧化剂的分类（掌握）：}
\begin{enumerate}[leftmargin=*]
    \item \textbf{脂溶性抗氧化剂：}
    \begin{itemize}[leftmargin=*]
        \item BHA（丁基羟基茴香醚）、BHT（二丁基羟基甲苯）、PG（没食子酸丙酯）、TBHQ（特丁基对苯二酚）。
        \item 主要用于油脂及含油食品，防止脂肪氧化酸败。
    \end{itemize}
    \item \textbf{水溶性抗氧化剂：}
    \begin{itemize}[leftmargin=*]
        \item 抗坏血酸（维生素C）及其钠盐、异抗坏血酸及其钠盐。
        \item 应用范围比脂溶性抗氧化剂更广。
    \end{itemize}
    \item \textbf{天然抗氧化剂：}
    \begin{itemize}[leftmargin=*]
        \item 茶多酚、维生素E（生育酚）、迷迭香提取物、竹叶提取物、甘草抗氧化物。
        \item 安全性优于合成抗氧化剂，日益受到重视。
    \end{itemize}
\end{enumerate}
\end{keybox}

\begin{tipbox}
\textbf{抗氧化剂的协同作用：}BHA与BHT混合使用时抗氧化效果优于单独使用。柠檬酸、磷酸等金属离子螯合剂可作为抗氧化剂的增效剂（Synergist），通过螯合催化氧化的金属离子来增强抗氧化效果。
\end{tipbox}

\subsection{着色剂（Color / Food Colorant）}

\begin{enumerate}[leftmargin=*]
    \item \textbf{天然色素：}
    \begin{itemize}[leftmargin=*]
        \item 辣椒红素（capsanthin）、姜黄素（curcumin）、红曲红（monascus red）、栀子黄、\(\beta\)-胡萝卜素、焦糖色。
        \item 优点：安全性较高，部分兼具营养功能（如\(\beta\)-胡萝卜素为维生素A前体）。
        \item 缺点：着色力较弱、稳定性较差（对光、热、pH敏感）、价格较高。
    \end{itemize}
    \item \textbf{合成色素：}
    \begin{itemize}[leftmargin=*]
        \item 柠檬黄（tartrazine）、苋菜红（amaranth）、亮蓝（brilliant blue）、日落黄、胭脂红、靛蓝。
        \item 优点：着色力强、色泽鲜艳、稳定性好、价格低廉。
        \item 缺点：安全性存疑（部分合成色素有致敏、行为影响等风险）。
    \end{itemize}
    \item \textbf{禁用色素（严禁用于食品！）：}
    \begin{itemize}[leftmargin=*]
        \item 苏丹红（Sudan Red）——致癌性，曾用于辣椒粉、鸭蛋（"红心鸭蛋"事件）
        \item 奶油黄（Butter Yellow）——致癌性
    \end{itemize}
\end{enumerate}

\subsection{防腐剂（Preservative）}

\begin{keybox}
\textbf{常用防腐剂分类（掌握）：}
\begin{enumerate}[leftmargin=*]
    \item \textbf{酸型防腐剂：}
    \begin{itemize}[leftmargin=*]
        \item 苯甲酸（Benzoic Acid）及其钠盐、山梨酸（Sorbic Acid）及其钾盐、丙酸（Propionic Acid）及其盐类。
        \item 有效形式为未解离的分子，因此\textbf{在酸性条件下（低pH）防腐效果更强}。在pH $>$ 5.5时效果显著下降。
        \item 山梨酸的安全性高于苯甲酸，是目前国际上应用最广泛的防腐剂之一。
    \end{itemize}
    \item \textbf{酯型防腐剂：}
    \begin{itemize}[leftmargin=*]
        \item 对羟基苯甲酸酯类（尼泊金酯类，Parabens）。
        \item 在较宽的pH范围内均有效（pH 4–8），不受酸性条件的限制。
        \item 缺点：水溶性差、具特殊气味。
    \end{itemize}
    \item \textbf{生物型防腐剂：}
    \begin{itemize}[leftmargin=*]
        \item \textbf{乳酸链球菌素（Nisin）：}由乳酸链球菌产生的一种多肽类抗菌物质。
        \item 特点：对G$^+$菌（特别是产芽孢菌如肉毒梭菌）有强烈抑制作用。对G$^-$菌、酵母、霉菌无效。
        \item 安全性高——可被人体消化道中的蛋白酶水解为氨基酸。
    \end{itemize}
\end{enumerate}
\end{keybox}

\subsection{护色剂（Color Fixative）— 亚硝酸盐}

\begin{keybox}
\textbf{护色剂（发色剂）——亚硝酸盐（Nitrite）*（掌握）：}
\end{keybox}

\begin{defbox}
\textbf{护色剂：}能使肉及肉制品呈良好色泽的物质。最常用的是\textbf{亚硝酸钠（NaNO$_2$）}和\textbf{硝酸钠（NaNO$_3$）}。
\end{defbox}

\begin{keybox}
\textbf{亚硝酸盐在肉制品中的三大作用*：}
\begin{enumerate}[leftmargin=*]
    \item \textbf{发色作用（护色）：}
    \begin{itemize}[leftmargin=*]
        \item 亚硝酸盐在肉中还原为NO $\rightarrow$ NO与肌红蛋白（Mb）反应生成亚硝基肌红蛋白（NOMb）$\rightarrow$ 加热后形成稳定的\textbf{亚硝基血色原}，赋予肉制品\textbf{鲜艳的桃红色}。
    \end{itemize}
    \item \textbf{抑菌作用（最重要的安全功能！）：}
    \begin{itemize}[leftmargin=*]
        \item 亚硝酸盐是目前已知\textbf{抑制肉毒梭菌（C.\ botulinum）最有效的物质！}
        \item 对肉毒梭菌芽孢的萌发和生长有强力抑制作用，是罐头类肉制品和腌制肉制品中\textbf{不可替代}的防腐成分。
    \end{itemize}
    \item \textbf{增强风味：}赋予腌腊肉制品特有的风味。
\end{enumerate}
\end{keybox}

\begin{exambox}
\textbf{亚硝酸盐危害的双面性（考试常见考点）：}
\begin{itemize}[leftmargin=*]
    \item \textbf{益处：}最佳的肉毒梭菌抑制剂 $+$ 良好的发色和风味作用。
    \item \textbf{风险：}与胺类物质在体内外结合生成$N$-亚硝基化合物（强烈致癌物）！尤其在胃的酸性条件下更易发生。
    \item \textbf{管理对策：}GB 2760规定亚硝酸盐在肉制品中的最大使用量（通常 $\leq$ 0.15 g/kg）及最终残留限量（$\leq$ 30 mg/kg），并要求与抗坏血酸或异抗坏血酸钠（VC可阻断亚硝化反应！）同时使用。
\end{itemize}
\end{exambox}

\subsection{甜味剂（Sweetener）}

\begin{enumerate}[leftmargin=*]
    \item \textbf{天然甜味剂：}
    \begin{itemize}[leftmargin=*]
        \item \textbf{糖醇类：}山梨糖醇（sorbitol）、木糖醇（xylitol）、麦芽糖醇（maltitol）、甘露糖醇（mannitol）。热量低、不引起血糖剧烈波动、不致龋齿。
        \item \textbf{非糖醇类：}甜菊糖苷（stevioside，从甜叶菊提取）、甘草酸。甜度极高（蔗糖的200–300倍）。
    \end{itemize}

    \item \textbf{人工合成甜味剂：}
    \begin{itemize}[leftmargin=*]
        \item \textbf{磺胺类：}甜蜜素（cyclamate）、糖精钠（saccharin）。历史悠久、价格低廉。甜蜜素在一些国家被禁用（膀胱致癌性存疑）。
        \item \textbf{二肽类：}阿斯巴甜（aspartame）——由天冬氨酸和苯丙氨酸组成。甜度为蔗糖的200倍！但苯丙酮尿症（PKU）患者禁用。热稳定性差，不宜用于高温加工食品。
        \item \textbf{蔗糖衍生物：}三氯蔗糖（sucralose）——以蔗糖为原料经氯代制成。甜度为蔗糖的600倍，不被人体代谢、无热量，耐高温。安赛蜜（acesulfame K）——甜度为蔗糖的200倍。
    \end{itemize}
\end{enumerate}

\begin{tipbox}
\textbf{阿斯巴甜的安全提醒：}必须在标签上注明"含苯丙氨酸"，以提醒苯丙酮尿症（PKU）患者避免食用。
\end{tipbox}

\subsection{增稠剂（Thickener）}

\begin{itemize}[leftmargin=*]
    \item \textbf{常用品种：}瓜尔胶、黄原胶、卡拉胶（carrageenan）、明胶、琼脂、羧甲基纤维素钠（CMC-Na）。
    \item \textbf{特征：}大多为大分子亲水性多糖类或多肽类，在肠道中不易消化吸收，大多数来自天然（植物、海藻、微生物发酵）。
    \item \textbf{功能：}提高食品黏度、改善口感和组织状态、稳定乳化体系、防止冰晶形成（冷冻食品中）。
\end{itemize}

\subsection{乳化剂（Emulsifier）}

\begin{itemize}[leftmargin=*]
    \item \textbf{常用品种：}卵磷脂（lecithin，从大豆或蛋黄提取）、单硬脂酸甘油酯（单甘酯）、聚甘油脂肪酸酯。
    \item \textbf{功能：}降低油-水界面张力，使不相溶的两相均匀分散，形成稳定的乳状液。面包中使脂肪均匀分布、延缓老化；冰淇淋中使脂肪均匀分散、改善口感。
\end{itemize}

\subsection{其他重要食品添加剂}

\begin{enumerate}[leftmargin=*]
    \item \textbf{漂白剂（Bleaching Agent）：}亚硫酸盐（还原性漂白）——用于干果、蜜饯、淀粉等。过量可破坏食品中维生素B1。GB 2760严格限制残留量（以SO$_2$计）。
    \item \textbf{膨松剂（Bulking Agent）：}碳酸氢钠、碳酸氢铵、明矾（硫酸铝钾——含铝，长期大量摄入有神经毒性风险）。
    \item \textbf{水分保持剂（Humectant）：}磷酸盐——提高肉制品保水性。
    \item \textbf{增味剂（Flavor Enhancer）：}味精（谷氨酸钠，MSG）、5$'$-呈味核苷酸二钠（IMP+GMP，与MSG有强烈协同增鲜效应）。
    \item \textbf{面粉处理剂：}过氧化苯甲酰（漂白面粉）——已在中国禁用（2011年起）。
\end{enumerate}

% =====================================================================
%  第五节 食品添加剂的管理制度
% =====================================================================
\section{食品添加剂的管理制度}

\begin{keybox}
\textbf{中国食品添加剂管理的法律依据：}
\begin{enumerate}[leftmargin=*]
    \item \textbf{《食品安全法》：}第39条规定国家对食品添加剂生产实行许可制度；第40条规定食品添加剂应当在技术上确有必要且经过风险评估证明安全可靠，方可列入允许使用的范围。
    \item \textbf{GB 2760《食品安全国家标准 食品添加剂使用标准》：}是食品添加剂使用的核心标准，规定了允许使用的添加剂品种、使用范围和最大使用量（或残留量）。
    \item \textbf{GB 14880《食品安全国家标准 食品营养强化剂使用标准》：}专门规定营养强化剂的使用范围和用量。
\end{enumerate}
\end{keybox}

\textbf{新食品添加剂的审批：}
\begin{itemize}[leftmargin=*]
    \item 审批机构：国家卫生健康委员会（NHCC，原卫生部）。
    \item 审批程序：申请人提交安全性评估资料、生产工艺、质量标准、使用范围等 $\rightarrow$ 食品安全风险评估中心评审 $\rightarrow$ NHCC批准公告 $\rightarrow$ 纳入GB 2760。
    \item 基本原则：技术上确有必要，且经风险评估证明安全可靠。
\end{itemize}

\textbf{食品添加剂的标识要求：}预包装食品的配料表中，食品添加剂应按GB 2760规定的通用名称标示；阿斯巴甜必须加注"含苯丙氨酸"字样。

% =====================================================================
%  复习题
% =====================================================================
\reviewcontent{
\section{复习题}

\subsection*{一、名词解释}

\begin{enumerate}[leftmargin=*, itemsep=6pt]
    \item \textbf{食品添加剂（中国《食品安全法》定义）、带入原则（Carry-over）}
    \item \textbf{ADI、GRAS}
    \item \textbf{护色剂、防腐剂、抗氧化剂}
\end{enumerate}

\subsection*{二、简答题}

\begin{enumerate}[leftmargin=*, itemsep=8pt]
    \item \textbf{【重点】}简述食品添加剂的定义及使用原则（基本要求）。
    \item \textbf{【重点】}简述JECFA对食品添加剂安全性评价的四类分级。
    \item \textbf{【重点】}简述亚硝酸盐在肉制品加工中的三大作用。为何亚硝酸盐虽然存在形成$N$-亚硝基化合物的风险，仍在肉制品加工中允许使用？如何控制这一风险？
    \item \textbf{【重点】}比较苯甲酸/山梨酸（酸型）与尼泊金酯类（酯型）防腐剂的作用特点差异。
    \item 简述抗氧化剂的分类，各举一例代表物质。
    \item 比较天然色素与合成色素的优缺点。
    \item 简述甜味剂的主要类型及阿斯巴甜使用的特殊注意事项。
    \item 简述食品添加剂"带入原则"的含义及其适用条件。
    \item 简述中国食品添加剂管理的主要法律法规和标准。
\end{enumerate}

\subsection*{三、综合论述题}

\begin{enumerate}[leftmargin=*, itemsep=8pt]
    \item \textbf{【综合】}结合GB 2760的使用原则，论述食品添加剂"安全性"与"技术必要性"的关系。试以亚硝酸盐为例，分析食品添加剂管理中"风险-获益"平衡的理念。
\end{enumerate}

\subsection*{参考答案要点}

\subsubsection*{名词解释（要点）}

\begin{enumerate}[leftmargin=*, itemsep=4pt]
    \item \textbf{食品添加剂：}为改善食品品质和色香味或防腐、保鲜和加工需要加入食品中的人工合成或天然物质。包括香料、胶基、加工助剂和营养强化剂。
    \item \textbf{带入原则：}添加剂非直接加入某食品，而是通过含此添加剂的配料带入。
    \item \textbf{ADI：}每日允许摄入量——终生每日摄入而无健康风险的剂量（mg/kg bw）。
    \item \textbf{GRAS：}第一类——一般公认安全，毒性极低，不规定ADI。
\end{enumerate}

\subsubsection*{简答题（要点）}

\begin{enumerate}[leftmargin=*, itemsep=6pt]
    \item \textbf{五项使用原则：}不产生健康危害、不掩盖腐败、不掩盖质量缺陷（不掺假）、不降低营养价值、最小使用量原则。
    \item \textbf{JECFA四类：}第一类GRAS（不设ADI）；第二类A（A1有正式ADI，A2有临时ADI）；第三类B（毒理资料不足，无ADI）；第四类C（C1不安全，C2仅特殊情况使用）。
    \item \textbf{亚硝酸盐三作用：}发色（形成亚硝基肌红蛋白）、抑制肉毒梭菌（最强抑菌剂）、增强风味。\textbf{风险-获益平衡：}其独特且不可替代的抑制肉毒梭菌的作用（可致死的食源性病原菌），使其在严格控制下利大于弊。\textbf{控制措施：}限量使用（$\leq$ 0.15 g/kg）、规定残留限量（$\leq$ 30 mg/kg）、与VC同时使用以阻断亚硝胺合成。
    \item \textbf{酸型防腐剂：}有效成分为未解离分子，低pH条件下效果好，pH $>$ 5.5效果剧降；\textbf{酯型（尼泊金酯）：}在pH 4–8均有效，不受酸性条件限制，但水溶性差。
    \item \textbf{抗氧化剂分类：}脂溶性（BHA/BHT/PG/TBHQ）、水溶性（抗坏血酸类）、天然（茶多酚/VE/迷迭香提取物）。
    \item \textbf{天然色素：}安全性高，但着色力弱、稳定性差、价贵。\textbf{合成色素：}着色力强、鲜艳、稳定、价廉，但安全性存疑。
    \item \textbf{甜味剂类型：}天然（糖醇类、甜菊糖苷）、人工合成（磺胺类—糖精、二肽类—阿斯巴甜、蔗糖衍生物—三氯蔗糖）。阿斯巴甜必须标注"含苯丙氨酸"以提醒PKU患者。
    \item \textbf{带入原则：}某食品中的添加剂由其原料带入，而非直接添加。需满足：原料允许使用、不超量、正常工艺、带入量显著低于直接添加水平。
    \item 法律：《食品安全法》第39/40条；核心标准：GB 2760（使用标准）和GB 14880（营养强化剂标准）；审批机构：国家卫健委（NHCC）。
\end{enumerate}

\newpage
}
